\documentclass[12pt]{article}
\usepackage[utf8]{inputenc}
\usepackage{amsmath}
\usepackage{amssymb}
\usepackage{graphicx}
\usepackage{hyperref}
\usepackage{geometry}
\geometry{a4paper, margin=1in}
\usepackage{xcolor}
\usepackage{listings}

\title{\textbf{LSTMs from Scratch \\ Part 2: Pure PyTorch Implementation}}
\author{Sovesh Mohapatra}
\date{\today}

\begin{document}
\maketitle

\section*{Introduction}
In Part 1, we explored the mathematics of gated recurrence---the forget gate, input gate, output gate, and the additive cell state update that creates a gradient highway through time. Today, we take that math and translate it into a pure PyTorch implementation. We will build from individual cells to complete sequence models.

\section*{The LSTM Cell}
Our LSTMCell class implements the four key computations. Each gate is a linear projection of the concatenated input $x_t$ and previous hidden state $h_{t-1}$, followed by a nonlinearity:

\begin{lstlisting}[language=Python, basicstyle=\small\ttfamily, commentstyle=\color{green!60!black}]
# Compute gates
i = sigmoid(W_xi @ x + W_hi @ h_prev)  # Input gate
f = sigmoid(W_xf @ x + W_hf @ h_prev)  # Forget gate
g = tanh(W_xc @ x + W_hc @ h_prev)     # Cell candidate
o = sigmoid(W_xo @ x + W_ho @ h_prev)  # Output gate

# Cell state update (the gradient highway)
c_new = f * c_prev + i * g

# Hidden state update
h_new = o * tanh(c_new)
\end{lstlisting}

\subsection*{Weight Initialization}
We use Xavier initialization for all weights, with a critical modification: \textbf{forget gate biases are initialized to 1.0}. This encourages the network to start with information preservation rather than information destruction, stabilizing early training. This trick, first noted by Jozefowicz et al. (2015), is essential for consistent convergence.

\begin{lstlisting}[language=Python, basicstyle=\small\ttfamily, commentstyle=\color{green!60!black}]
# Critical: Initialize forget gate bias to 1.0
nn.init.ones_(self.Wxf.bias)
nn.init.ones_(self.Whf.bias)
\end{lstlisting}

\section*{Multi-Layer LSTMs}
Stacking LSTM cells creates hierarchical temporal processing:
\begin{itemize}
    \item Lower layers capture short-term patterns (local token interactions)
    \item Higher layers learn long-term structure (global sentence-level meaning)
    \item Dropout between layers prevents overfitting
\end{itemize}

The output of each layer's hidden state $h_t^{(l)}$ becomes the input for the next layer. We apply dropout only between layers, never at the final output---matching the convention used in PyTorch's built-in LSTM.

\section*{Architecture Variants}

\subsection*{LSTM Classifier (Many-to-One)}
For sequence classification tasks:
\begin{itemize}
    \item Processes the entire sequence through multi-layer LSTM
    \item Concatenates final hidden states from \textit{all layers}
    \item Projects to class logits via a fully-connected layer
    \item Optional bidirectional support doubles the representation power
\end{itemize}

\subsection*{LSTM Tagger (Many-to-Many)}
For sequence labeling (POS tagging, NER):
\begin{itemize}
    \item Outputs a prediction at \textit{each} time step
    \item Projects each hidden state $h_t$ to output vocabulary
\end{itemize}

\subsection*{CharLSTM}
Character-level language model:
\begin{itemize}
    \item Embedding layer maps character indices to dense vectors
    \item Multi-layer LSTM processes the embedded sequence
    \item Autoregressive generation uses temperature-controlled sampling
\end{itemize}

\begin{lstlisting}[language=Python, basicstyle=\small\ttfamily, commentstyle=\color{green!60!black}]
def generate(self, start_token, seq_len, temperature=1.0):
    # Sample next character from the distribution
    next_logits = logits[:, -1, :] / temperature
    probs = softmax(next_logits, dim=-1)
    next_token = torch.multinomial(probs, 1)
\end{lstlisting}

\subsection*{Seq2SeqLSTM (Encoder-Decoder)}
The classic encoder-decoder architecture:
\begin{itemize}
    \item \textbf{Encoder} LSTM compresses the source sequence into a fixed-size context vector (its final hidden and cell states)
    \item \textbf{Decoder} LSTM is initialized with the encoder's final states and generates the target sequence autoregressively
\end{itemize}

\section*{Bidirectional Processing}
For tasks requiring full context (e.g., sentiment analysis, NER), we implement bidirectional LSTMs:
\begin{itemize}
    \item Forward LSTM processes left-to-right: $h_t^{\rightarrow}$
    \item Backward LSTM processes right-to-left: $h_t^{\leftarrow}$
    \item Concatenate both: $h_t = [h_t^{\rightarrow}; h_t^{\leftarrow}]$
\end{itemize}

\section*{Training Strategy}
We train on a challenging long-range dependency task specifically designed to expose the vanishing gradient problem:
\begin{itemize}
    \item Input: random sequence of length 30
    \item Task: classify based on the \textit{first + last} element only
    \item The model must remember information across 30 time steps while ignoring irrelevant intermediate values
\end{itemize}

Gradient clipping prevents the complementary exploding gradient problem:
\begin{lstlisting}[language=Python, basicstyle=\small\ttfamily, commentstyle=\color{green!60!black}]
torch.nn.utils.clip_grad_norm_(model.parameters(), max_norm=1.0)
\end{lstlisting}

\section*{Next Steps: Visualizing Gate Dynamics}
With our LSTM implemented and architecture variants in place, we can now train and analyze \textit{what the gates actually learn}. In \textbf{Part 3}, we will visualize forget gate, input gate, and output gate activations across time---watching the network learn to selectively remember and forget.

\vspace{1em}
\noindent \textit{Full code is on the GitHub repo. Stay tuned for the gate visualization drop!}

\end{document}
